%% 01-introduction.tex --- included by ../sfc.tex
\section{Introduction}
\label{sec:intro}

A distributed quantum computer (DQC) runs one circuit across several small
cores joined by links that distribute entanglement on
demand~\cite{cuomo2020towards,escofet2024review}. Any two-qubit gate
whose operands sit in different cores must be paid for with
Einstein--Podolsky--Rosen (EPR) pairs, and distributing a pair is far
slower and less reliable than any operation inside a
core~\cite{stephenson2020high,daiss2021quantum}. Pairs consumed is
therefore the metric we minimise (Section~\ref{sec:scope}).

A pair can be spent three ways. Teleporting a \emph{state} (\TeleData{})
relocates a logical qubit, making every later gate it meets in its new
core local~\cite{bennett1993teleporting}. Teleporting a \emph{gate}
(\TeleGate{}) executes one remote CX through a cat-entangler with neither
operand moving~\cite{eisert2000optimal,yimsiriwattana2004generalized}.
Holding one cat-entangled resource open across a run of remote CX gates
sharing a control costs one pair for the whole run --- a \emph{packet},
following AutoComm's gate-sequence
amortisation~\cite{wu2022autocomm}. Which is cheapest depends on the
circuit structure where it is applied, so none dominates.

The choice must be made before its price is known. Each alternative
becomes executable only after its own preparation: the right operand
beside the right communication port, that port free, the destination
holding a spare slot. A compiler performs that work only for the
alternative it selects. Scoring an action together with the preparation
it currently needs therefore ranks candidates by the layout the compiler
happens to occupy rather than by the progress they buy, and staging a
losing candidate still perturbs the layout.

\PORTER{} --- a \emph{Plan-Oriented Router for Teleportation-Enabled
aRchitectures} --- separates the two questions. Its unit of choice is a
\emph{plan}: a complete specification of one inter-core action and the
conditions it needs, omitting the SWAP sequence that would establish
them. Plans are compared on the gate progress they buy, one is committed,
and only then are preparation SWAPs generated for it. Around this loop
\PORTER{} adds two-tier capacity relief, which manufactures free slots
when every candidate is blocked, and a deterministic seed-free initial
layout in place of the seeded best-of-$N$ search its predecessors run.

\textbf{Results.} We evaluate on MQT-Bench~\cite{quetschlich2023mqtbench}
suites at 25, 64 and 80 logical qubits, on a $100$--$360$-qubit
scale-up series, and on two networks of IBM heavy-hex cores --- forty-eight
circuit--architecture instances. Against
\TeleSABRE{}~\cite{russo2025telesabre}, the published state of the art
for this cost model, \PORTER{} uses substantially fewer EPR pairs: over
the instances both complete it lowers geometric-mean consumption by
$59.6\%$, $41.6\%$ and $68.2\%$ on the three suites at the same
three-attempt budget, and it completes every instance,
including the seven --- among them an 80-qubit, $23{,}381$-CX circuit
--- on which \TeleSABRE{} does not converge in three seeds.

Against \dSABRE{}~\cite{lidsabre} the picture is mixed and we report
every point of it. Under each compiler's own protocol \PORTER{} costs
$1.07\times$ its pooled geometric-mean EPR at $0.58\times$ its compile
time; a comparable pass budget reaches $1.05\times$ at $0.74\times$. From
identical externally supplied layouts at matched budgets \PORTER{} leads
across forty-one instances ($126.6$ against $134.2$), and the same
control takes both heavy-hex core graphs. On the $100$--$360$-qubit
series it leads outright, at $0.72\times$, which one mistuned constant
explains: the lookahead horizon (Section~\ref{sec:horizon}). The two
compilers are close relatives and share four design decisions, set out in
Section~\ref{sec:related} so neither is credited with the other's.

\textbf{What the mechanisms are worth.} A head-to-head number says little
about which design decision earns its place, so we ablate each in
isolation. Capacity relief is what makes a deterministic-greedy design
complete in practice: without it four of six 80-qubit instances fail. The
initial layout separates the two compilers below roughly a hundred
physical qubits; the lookahead horizon separates them above it. Two
answers are negative. Plan \emph{commitment} --- the rule that a selected
plan is not reopened until ready --- reaches the same EPR as re-scoring
the front every iteration, within $0.6\%$ and on the same number of
moves, so it buys compile time rather than better routes. That ablates
the commitment rule, not the plan abstraction beneath it, which no
experiment here removes (Section~\ref{sec:discussion}). And the packet
primitive is beaten at 80 qubits by removing gate teleportation
altogether, by $32.8\%$ --- an interaction between action vocabulary and
plan scorer, not a verdict on amortisation.

\textbf{Problem.} Given a circuit DAG on $n$ logical qubits, a machine of
$K$ cores with intra-core coupling graphs, inter-core links and designated
communication ports, and an initial layout, emit a legal operation stream
realising every gate while minimising the number of EPR pairs consumed.
Section~\ref{sec:scope} states the cost model this objective assumes and
what it excludes.

\textbf{Contributions.} We make three.

\begin{enumerate}[leftmargin=1.5em,itemsep=2pt,topsep=2pt,label=(\roman*)]
\item \emph{The \PORTER{} compiler}: a plan-oriented remote stage that
  scores a complete inter-core action before any preparation SWAP for it
  exists and commits before preparing
  (Section~\ref{sec:sfc-plan-first}), a capacity discipline that treats
  free slots as a legality condition and manufactures them by a two-tier
  relief search (Section~\ref{sec:sfc-relief}), and a deterministic
  seed-free initial layout replacing the seeded best-of-$N$ search its
  predecessors run (Section~\ref{sec:sfc-layout}).
\item \emph{The lookahead horizon has no good constant.} Every
  \SABRE{}-descended router fixes it by constant. Its best value spans a
  factor of thirty-two across our benchmark, and choosing it per circuit
  by routing a small pruned candidate set is what carries \PORTER{} past
  \dSABRE{} at scale (Section~\ref{sec:horizon}).
\item \emph{A controlled study over forty-eight instances}, under
  native, shared-layout and cross-layout protocols, with ablations that
  attribute the result to individual mechanisms and report two negative
  findings (Sections~\ref{sec:results}--\ref{sec:discussion}). Every
  route is certified by a separately implemented replay verifier, and
  the harness reproducing every table is released.
\end{enumerate}

Section~\ref{sec:rqs} states the four questions the evaluation answers.
Each mechanism is tied to the experiment that assesses it: plan
commitment to the hop-wise ablation, capacity relief to the completion
ablation, the initial layout to the shared- and cross-layout controls,
horizon selection to the scale-up series and the star-Multiplier
livelock, packets to the action-vocabulary study, and the verifier to
fault injection.

The \PORTER{} implementation, the replay verifier, the full cross-router
comparison harness, and the raw per-circuit results are publicly available at
\url{https://github.com/ebony72/porter}.

%%────────────────────────────────────────────────────────────────────────────
